% Transcription — fonds Alexandre Grothendieck, shelfmark 80, batch 2 (pages 21-40).
% Established from the facsimile archives/batches/80/batch-02.pdf, cut from a
% copy of 80.pdf fetched through the project's own relay
% (https://grothendieck-relay.onrender.com/source/80.pdf), not uploaded by the
% user; PDF page k = archivists' page 20+k, no cover sheet. Per the skill
% transcribe-grothendieck. The folder is seventy-two pages in four batches
% (1-20, 21-40, 41-60, 61-72), transcribed in parallel; this is the second.
% Batch 1 (transcripts/80/batch-01.fr.tex) was read at the end of this pass
% and its notation followed (pseudo-droites abbreviated « ps. droites » as he
% writes it, \mathfrak{S} for his gothic S, configurations numbered as he
% numbers them); it was not edited.
%
% Pass: Opus 5.5 (claude-opus-5-5), 2026-09-26 — first pass, unchecked against the pages by a human.
%
% Dating and title. \dating is the folder's own date, copied verbatim from
% src/content/catalogue.ts; the group « Géométrie et topologie combinatoire »
% (69-89) is not added, since the folder has a date of its own. Title from the
% catalogue, without its final full stop, as in batch 1 of this folder.
% Page 27 is a typed USTL notice dated « Montpellier, le 13 mai
% 1981 », which is consistent with the archivists' « [à partir de 1981] »; it
% carries nothing of his and is skipped.
%
% Pages skipped: 21, 23, 25 and 29, photocopies of a handwritten letter in
% German, autobiographical and unrelated to the mathematics (two copies each of
% two leaves); 27, the typed administrative notice above; 30, a blank sheet.
% Pages summarised: none. Pages transcribed from his hand: 22, 24, 26, 28
% (black ink, then brown ink, fast drafting) and 31-40 (brown ink, a fair run).
%
% His pagination: pages 22-28 carry none. Pages 31-40 are one run numbered on
% the rectos only, top left: 1 (circled, p. 31), 2 (p. 33), 3 (p. 35),
% 4 (p. 37), 5 (p. 39); the even pages are the versos and continue the text.
% Recorded in a \note at each numbered page. His own equation numbers (1)-(18),
% (6 bis), (11 bis), (15 bis) run through pp. 31-39 and are kept as he writes
% them.
%
% What the batch is about.
%   22-28  Continuing batch 1's p. 20: a classification of the « crucial »
%          1-dimensional arrangements of N pseudo-lines, by types he names
%          I_N (a pencil of N-1 lines plus a transversal), I_{p,q} (two pencils
%          and the line through their centres), II'_{p,q} (two pencils, no
%          joining line), II_4, II_5 (simple arrangements of 4 and 5 lines),
%          III_6 (tetrahedral arrangement), IV_7 (complete quadrilateral with
%          its three diagonals), and a type I'_{3,2} / II_{3,2} that is not
%          defined in these pages. P. 24: « Théorème », every arrangement
%          without bigons is II_4 or contains one of I'_{3,2}, II_5, III_6
%          (boxed), and these are « latticiel »: closed facets form a lattice.
%          P. 26: « Cas latticiel », characterised by no bigonal line and no
%          II_4; a Proposition on adding an (N+1)-th non-bigonal line D (two
%          vertices s, s' on D adjacent iff a closed facet contains both); a
%          struck Lemme. P. 28: the Lemme again, a numbering D_1..D_N with
%          D_1..D_4 of type II_4 (resp. D_1..D_6 of type III_6) and D_i
%          non-bigonal in K_i for i >= 5, with an « Idée » of proof by cases.
%   31-40  « I) Arrangements de 3 pseudo-droites » (his heading): the
%          combinatorial cube and octahedron from three 2-element sets Phi_i,
%          faces/edges/vertices as partial sections, the antipody a, quotients
%          by a (« gauche » or projective cube and octahedron: 3 vertices, 6
%          edges, 4 faces, i.e. three lines in RP^2); an equivalence of
%          oriented objects with « gauche » ones (5); orientations via the
%          « exterior product » of 2-element sets, or(C) ≃ (∧ Phi_i) ∧ or(I)
%          (6), edge sets A_i, vertex sets S_i ≃ I \ {i}, orientation sets
%          Omega_i = S_i ∧ A_i (9)-(15 bis); the canonical trivialisation of
%          the Z/2Z-torsor (15 bis) coming from the map, hence the equivalence
%          (16) of « oct. gauches » with « graphes octaédraux gauches » with a
%          trivialisation ∧ A_i ≃ 1; the commutative cube (17) of octahedra /
%          octahedral graphs, gauche / ordinary, oriented / not, commented on
%          p. 38; p. 39 defines an orientation of a 1-dimensional arrangement
%          (K, (D_i)_{i in J}) as orientations omega_I of the octahedral graphs
%          of its triples, the functor (18) X -> (K, (D_i), omega), and a
%          « Th. »: (18) is fully faithful; p. 40 starts over with the axioms
%          of a 1-dimensional arrangement of pseudo-lines.
%
% Notation fixed once for the batch: roman numerals of the types as
% \mathrm{I}, \mathrm{II}, ... (he draws them with serifs and sometimes an arc
% over them); \wedge / \bigwedge for his « exterior product » of 2-element
% sets; \underline{1} for his one-element set; \mathfrak{P}_k for his gothic
% P (sets of k-element subsets); \mathcal{X} for his cursive barred X (a map /
% arrangement); \underline{a} for the antipody (underlined on the page).
%
% Readings to check first: the name I'_{3,2} / II_{3,2} on pp. 24-28 and its
% definition, which is not in this batch; the prose of pp. 26 and 28 (many
% \ill{}); p. 38, a fast commentary whose syntax is partly reconstructed; the
% marginal notes on pp. 33, 35, 37, 38, 39, 40; the arrow labels of the cube
% (17) and the small hook marks on its vertical arrows.

\documentclass[11pt,a4paper]{article}
% Le préambule commun à toutes les transcriptions du fonds Grothendieck.
%
% Il tient en une idée : **l'incertitude se compose, elle ne se résout pas.**
% Une transcription qui lisse un mot douteux a détruit la seule chose pour
% laquelle on la faisait — un lecteur doit pouvoir savoir, à chaque endroit,
% ce qui était sur le feuillet et ce qui a été ajouté. Les six macros qui
% suivent sont donc l'appareil critique, et non de la décoration.
%
% Les mêmes commandes sont reconnues par `scripts/render.mjs`, qui produit la
% vue de lecture du site. Une macro ajoutée ici doit l'être aussi là-bas,
% faute de quoi le rendu échouera bruyamment — ce qui est voulu : un rendu
% silencieusement faux est pire qu'un rendu absent.

\ProvidesPackage{grothendieck}[2026/08/08 Transcriptions du fonds Grothendieck]

\usepackage{amsmath,amssymb,amsthm}
\usepackage{mathrsfs}
\usepackage{stmaryrd}
% Les diagrammes commutatifs sont partout dans ce fonds : c'est la langue
% même de la géométrie algébrique de Grothendieck, et une transcription qui
% les rendrait en prose ne transcrirait rien.
\usepackage{tikz-cd}
% Français par défaut — c'est la langue des feuillets — mais l'anglais est
% chargé aussi : la traduction et le résumé sont des documents anglais, et la
% typographie française (espaces avant les deux-points) y serait une faute.
% Ces documents basculent par \selectlanguage{english} en tête de corps.
\usepackage[english,french]{babel}
\usepackage{fontspec}
\usepackage{geometry}
\geometry{a4paper,margin=25mm,marginparwidth=28mm}
\usepackage{marginnote}
\usepackage{xcolor}
% ulem, not soul. `soul`'s \st and \ul are built on a character-by-character
% scan that XeLaTeX's font machinery breaks: the document fails with an
% undefined control sequence rather than degrading. ulem does the same two jobs
% and is Unicode-engine safe. `normalem` stops it hijacking \emph, which the
% transcriptions use for Grothendieck's own emphasis.
\usepackage[normalem]{ulem}
\usepackage{hyperref}
\hypersetup{colorlinks=true,linkcolor=black,urlcolor=black,pdfborder={0 0 0}}

% --- Métadonnées du lot -----------------------------------------------------
% Reprises telles quelles par le script de rendu, qui les affiche en tête de
% la vue de lecture. Elles fixent ce que le fichier prétend couvrir : sans
% elles, une transcription flotte sans point d'attache dans un fonds de seize
% mille feuillets.

\newcommand{\folder}[1]{\gdef\grothFolder{#1}}
\newcommand{\batch}[1]{\gdef\grothBatch{#1}}
\newcommand{\leaves}[2]{\gdef\grothFirst{#1}\gdef\grothLast{#2}}
\newcommand{\foldertitle}[1]{\gdef\grothFoldertitle{#1}}
\gdef\grothFolder{?}\gdef\grothBatch{?}\gdef\grothFirst{?}\gdef\grothLast{?}\gdef\grothFoldertitle{}

% --- Le résumé --------------------------------------------------------------
%
% Il ouvre la lecture modernisée, sous le titre « Résumé », et il est composé
% en petit. Ce n'est pas un ornement : c'est le seul passage du document qui
% ne soit pas des mathématiques, et un lecteur doit pouvoir le reconnaître
% avant de l'avoir lu — puis le sauter s'il connaît déjà le sujet, ou s'y
% arrêter s'il ne le connaît pas. Le filet qui le suit marque la reprise du
% corps.

\newenvironment{resume}
  {\small\setlength{\parskip}{0.45em}}
  {\par\medskip\hrule\medskip}

% --- Mots-clés ---------------------------------------------------------------
%
% En anglais, en clôture du résumé de la lecture modernisée : le vocabulaire
% moderne sous lequel chercher ce que les feuillets construisent. Le manifeste
% du site (`npm run manifest`) les extrait pour étiqueter la cote entière —
% cette ligne est la source unique des tags, il n'y a pas de fichier à côté.
\newcommand{\keywords}[1]{\par\smallskip\noindent
  {\footnotesize\textsc{Keywords} --- \textit{#1}}\par}

% --- L'appareil critique ----------------------------------------------------

% Le numéro de feuillet, en marge. C'est l'ancre qui permet de régler un
% désaccord sur une formule en regardant *un* feuillet plutôt qu'une chemise
% de six cents. Le site s'en sert aussi pour tourner les pages du fac-similé
% quand on fait défiler la transcription.
\newcommand{\leaf}[1]{%
  \marginnote{\footnotesize\color{gray!70}\textsf{#1}}%
  \label{leaf:#1}\ignorespaces}

% Illisible : on ne devine pas. Un mot inventé qui se lit comme les autres est
% le pire résultat possible.
\newcommand{\ill}{\textcolor{red!60!black}{[\,\ldots\,]}}

% Lecture douteuse : le mot est donné, et signalé comme incertain.
\newcommand{\uncertain}[1]{\uline{#1}}

% Ajout de l'éditeur — une lettre manquante, un mot sous-entendu.
\newcommand{\add}[1]{\textcolor{blue!50!black}{[#1]}}

% Biffé par Grothendieck lui-même. Ce qu'il a rayé fait partie du document :
% une rature dit souvent où la pensée a changé de direction.
\newcommand{\struck}[1]{\sout{#1}}

% Note du transcripteur, sur l'état matériel du feuillet.
\newcommand{\note}[1]{\par\smallskip{\small\color{gray!60!black}\textit{[#1]}}\par\smallskip}

% Note marginale de Grothendieck, distincte du corps du texte.
\newcommand{\marginal}[1]{\par\smallskip\noindent\hspace{1em}%
  {\small\color{gray!60!black}#1}\par\smallskip}

% --- Titre ------------------------------------------------------------------

\AtBeginDocument{%
  \begin{center}
    {\large\bfseries Fonds Alexandre Grothendieck --- cote n\textsuperscript{o}~\grothFolder}\\[2pt]
    {\itshape\grothFoldertitle}\\[4pt]
    {\small feuillets \grothFirst--\grothLast\ (lot \grothBatch)}\\[2pt]
    {\footnotesize\color{gray!60!black}Transcription établie d'après le fac-similé numérisé
     par l'Université de Montpellier.\\
     Les lectures incertaines sont soulignées, les illisibles notées [\,\ldots\,],
     les ajouts entre crochets.}
  \end{center}
  \vspace{1em}}

\endinput


\folder{80}
\batch{2}
\pages{21}{40}
\dating{[à partir de 1981]}
\watermark{Édition de démonstration}
\foldertitle{Systèmes de pseudo-droites : notes manuscrites (s.d.)}

\begin{document}

\page{22}
Cas cruciaux d'arrangement $1$-dimensionnels de $N$ ps.\ droites, avec $N \geq{}$ \struck{$5$} $4$ (le cas $N = 3$ étant \struck{\ill{}} et bien étudié)
\note{les chiffres romains des types ($\mathrm{I}$, $\mathrm{II}$, \ldots) sont tracés avec leurs barres horizontales en haut et en bas, parfois prolongées en arc au-dessus du nombre ; on les compose en romain droit, sans surligner}

1°) Cas où il existe une ps.\ droite bigonale.
\note{à droite, un dessin : un faisceau de droites concourantes en un point, coupé par une droite transversale}

$\mathrm{I}_N$) faisceau de $n = N-1$ ($\geq 3$) ps.\ droites \struck{passant par un \ill{}} \add{\uncertain{concourantes}} et d'une ps.\ droite \uncertain{sécante} (cas où la ps.\ droite bigonale \ill{})
\marginal{$N \geq 4$, $N+1$ sommets}
\note{« concourantes » est écrit au-dessus de la fin de ligne biffée ; la parenthèse finale est écrite en interligne, un mot y est biffé avant « ps.\ droite »}

$\mathrm{I}_{p,q}$) Système formé de deux faisceaux de $p$ et $q$ ps.\ droites concourantes ($p \geq q \geq 2$) plus \struck{\ill{}} une ps.\ droite passant par les $2$ sommets des deux faisceaux (\uncertain{Cas de} ps.\ droite bigonale \uncertain{unique}) $(p+q+1 = N)$
\marginal{$pq+2$ sommets ; $N \geq 5$}
\note{à gauche, sous « $N \geq 5$ », un dessin : deux faisceaux de droites et la droite qui joint leurs deux sommets}

2°) Il n'existe pas de ps.\ droite bigonale.

$\mathrm{II}'_{p,q}$) Système formé de deux faisceaux de $p$ et $q$ ps.\ droites concourantes, $p \geq q \geq 2$, $p+q = N \geq 5$ (les deux faisceaux \uncertain{i.e.}\ les \ill{} sommets \uncertain{des} \ill{} \ill{})
\marginal{$pq+2$ sommets ; $N \geq 5$}
\note{le symbole est d'abord écrit $\mathrm{II}$ puis surchargé, le prime ajouté ; il est lu $\mathrm{II}'_{p,q}$ d'après la page 24. À gauche, un dessin : deux faisceaux dont les sommets ne sont pas joints, la droite qui les joindrait en pointillé}

Le cas « \uncertain{minimal} » $p ={}$ \struck{$3$} $2$, $q = 2$ est \uncertain{justiciable} \ill{} \ill{}.

$\mathrm{II}_4$) \struck{Arrangement \uncertain{général} de $4$ ps.\ droites en position} Arrangement simple de $4$ ps.\ droites
\marginal{$6$ sommets}

$\mathrm{II}_5$) Arrangement simple de $5$ ps.\ droites.
\marginal{$10$ sommets}

\struck{Dans le cas où il n'y a pas de ps.\ droite bigonale (i.e.\ pas $\mathrm{I}_N$ ni $\mathrm{I}_{p,q}$), \ill{}}

\page{24}
\emph{Théorème} : supposons qu'il n'existe pas deux sommets tels que \struck{\ill{}} les ps.\ droites passent par l'un d'eux, i.e.\ qu'on ne soit pas dans l'un des cas $\mathrm{I}_N$, $\mathrm{I}_{p,q}$, $\mathrm{II}'_{p,q}$. Supposons de plus \struck{$N \geq 5$ (i.e.\ \ill{} \ill{} le cas $\mathrm{II}_4$)} \struck{et} que $K$ ne contienne pas de \uncertain{sous-arrangement} du type $\mathrm{II}_5$ ; elle est de l'un des types suivants : \struck{\ill{}}
\note{« Théorème » est souligné ; les deux points qui suivent « suivants » précèdent un signe biffé}

$\mathrm{III}_6$) Arrangement tétraédral
\marginal{$4+3 = 7$ sommets}
\note{à droite, un dessin : six droites, les côtés d'un triangle et trois droites issues de ses sommets concourant en un point intérieur ; quatre points triples marqués}

$\mathrm{IV}_7$) Arrangement d'un quadrilatère \uncertain{complet} et de ses trois diagonales
\marginal{$6+3 = 9$ sommets}
\note{au-dessous, un dessin : sept droites, dont certaines en tirets ; points marqués en noir, deux points entourés d'un cercle}

\emph{NB} $\mathrm{III}_6$ et $\mathrm{IV}_7$ \uncertain{sont} (\uncertain{les} \struck{\ill{}} \ldots \ldots $\mathrm{I}_{2,2}$ (\uncertain{quadrilatère} \uncertain{avec} $1$ \uncertain{diagonale}) et $\mathrm{I}_N$ (\uncertain{NB})
\note{fin de ligne très reprise : « $\mathrm{I}_{2,2}$ (\ldots) et $\mathrm{I}_N$ » est écrit au-dessus, et relié par un trait, d'une ligne biffée où l'on devine « \uncertain{configurations} » ; l'ensemble n'est pas sûr}

\ldots $N \geq 5$ qui ne contiennent pas l'arrangement $\mathrm{I}'_{3,2}$ ni $\mathrm{II}_5$. \uncertain{Tout} arrangement sans bigones, \struck{\ill{} types et qui \ill{}} est soit du type $\mathrm{II}_4$, soit contient un des types
\[
\boxed{\mathrm{I}'_{3,2},\ \mathrm{II}_5,\ \mathrm{III}_6}
\]
(c'est s'il ne contient $\mathrm{I}'_{3,2}$ ni $\mathrm{II}_5$, c'est un type $\mathrm{III}_6$ \uncertain{ou} $\mathrm{IV}_7$). Ces trois types \struck{sont} « \uncertain{latticiels} », i.e.\ dans le cas $2$-dimensionnel, les \uncertain{facettes} fermées forment un \uncertain{lattice} : \struck{\ill{}} l'intersection de deux \uncertain{facettes} fermées non vides est une \uncertain{facette}.
\note{le nom $\mathrm{I}'_{3,2}$ n'a pas été défini sur les pages précédentes de ce lot ; le prime et les indices $3,2$ sont nets. La liste encadrée est entourée d'un cadre arrondi. Comparer la page 20 (lot 1), où la même propriété est dite « non spéciales » : l'intersection de deux faces fermées est vide, réduite à un sommet, ou est une arête fermée}

\page{26}
\emph{Cas latticiel} : l'intersection de deux \uncertain{facettes} fermées est vide, ou une facette fermée.
\note{la page est à l'encre brune}

\emph{Prop.} Pour qu'on soit dans le cas latticiel, il faut et suffit que a) Pas de ps.\ droites \uncertain{bigonales}, et b) pas de type $\mathrm{II}_4$.

\emph{Proposition} Soit $\mathcal{X}$ un arrangement $2$-dim.\ de $N$ ps.\ droites $(D_i)_{i \in I}$, et considérons une $(N+1)$\add{ième} ps.\ droite $D$ d'un $\mathcal{X}'$ qui ne soit pas bigonale de $\mathcal{X}'$ ; soit $S$ l'ens.\ des sommets sur $D$, i.e.\ l'ens.\ des intersections de $D$ avec les $D_i$ $(i \in I)$. \struck{Pour \ill{}} \struck{pour $(D_i)$} \struck{$s \in S$, soit $\tilde{s}$ la facette \ill{}.} Pour que $s, s' \in S$ soient adjacents sur $D$, il faut et suffit qu'il existe une facette fermée $f$ de $\mathcal{X}$ qui contienne $s$ et $s'$ --- celle-ci est alors unique, dite \uncertain{intermédiaire}, et $D \cap f$ est un \ill{} \ldots\ $s$, $s'$, tel que $T \cap f = \lbrace s, s' \rbrace_T$, \ill{} sur $D$ \uncertain{joignant} $s$ et $s'$.
\note{« $(D_i)_{i \in I}$ », « $(N+1)$ième », « d'un $\mathcal{X}'$ » et « de $\mathcal{X}'$ » sont des ajouts interlinéaires ; $\mathcal{X}$ rend son X cursif barré. La fin de phrase, autour de $D \cap f$ et de $T \cap f$, est d'une lecture incertaine ; l'indice $T$ après l'accolade est tel quel}
\note{à gauche, un dessin : quelques droites formant une facette ; sur deux d'entre elles les sommets $s$ et $s'$, joints par une courbe marquée $T$ ; une droite marquée $J$ en bas}

Segment \uncertain{joignant} \struck{\ill{}} c'est aussi l'\uncertain{arête} sur $D$ joignant $s$ et $s'$.

\struck{\emph{Lemme} Pour un arrangement $1$-dim.\ de $N$ ps.\ droites \struck{$N \geq 5$} \ill{} sans bigones, \ill{} $\mathrm{II}_4$, et un ps.\ de \ldots\ \ill{} \ill{} \ill{} des types $\mathrm{II}_5$, $\mathrm{II}_{3,2}$ ou $\mathrm{III}_6$. Quand on rajoute des ps.\ droites, elles ne sont jamais bigonales.}
\marginal{$\mathrm{II}_4$, $\mathrm{II}_5$, $\mathrm{II}_{3,2}$, $\mathrm{III}_6$}
\note{tout ce paragraphe est barré de trois longues diagonales, et plusieurs lignes y sont en outre biffées ; il est repris à la page 28. La liste écrite à gauche, sous un trait, porte d'abord $\mathrm{II}_{3,1}$ (?) corrigé en $\mathrm{II}_{3,2}$ ; au-dessous, barré, « \ill{} \ill{} \ill{} repoussent des ps.\ droites »}

\page{28}
\emph{Lemme} Soit \struck{$\mathcal{X}$} $K$ un arrangement $1$-dimensionnel de $N$ ps.\ droites, sans ps.\ droite bigonale, des \ill{} $N \geq 5$. Alors on peut numéroter \struck{\ill{}} en $D_1, \ldots, D_N$, de telle façon que, si $K_i$ désigne l'arrangement défini par $D_1, \ldots, D_i$, donc, pour \struck{\ill{}} (\uncertain{resp.}\ $D_1, \ldots, D_6$ type $\mathrm{III}_6$)
\begin{itemize}
\item[1°)] $D_1, \ldots, D_4$ de type $\mathrm{II}_4$ (\uncertain{ou} \ill{} $(i \geq 7)$)
\item[2°)] Pour $i \geq 5$, $D_i$ non bigonale dans $K_i$. \struck{pour $i \geq 7$, $D_i$ \ill{}} (\emph{NB} Dans le cas \uncertain{inverse}, \ill{} dans le \struck{\ill{}} \ill{})
\end{itemize}
\note{le « $1$-dimensionnel » est écrit au-dessus de la ligne ; « (resp.\ $D_1, \ldots, D_6$ type $\mathrm{III}_6$) » est ajouté à droite de « $D_1, \ldots, D_4$ », « \uncertain{ou} \ill{} $(i \geq 7)$ » sous « de type $\mathrm{II}_4$ ». La parenthèse du NB n'est pas fermée ; un mot de quatre lettres y est biffé, lu « UPS » sans certitude}

\emph{Idée} 2°) \ill{} \ill{} \uncertain{automatique}. \ill{} \struck{\ill{}} \ill{}, \ill{} que $K_5$ ne soit pas \ill{}, i.e.\ $K_5$ est du type $\mathrm{II}_5$ ou du type $\mathrm{II}_{3,2}$ ; \ill{} $K_i$ ($i \geq 6$) \ill{} \ill{} \ill{} $\mathrm{II}_{\ill{}}$, et dans les autres types \ill{} \ill{} $K$ \ill{} \ill{} et que \ill{} $\mathrm{III}_6$ \ill{} \ill{} $\mathrm{III}_7$. Dans le même plan \ill{} au cas $\mathrm{III}_6$, \ill{} \ldots
\note{passage très rapide, dont seules les formules se lisent sûrement}

On sait $\exists\, K_5$ du type $\mathrm{II}_5$ --- et quand on continue à rajouter des ps.\ droites \ill{} \ill{} $1 : 1$, c'est par des bigones.

On sait $\exists\, K_5$ du type $\mathrm{II}_{3,2}$, et $K$ de type $\mathrm{I}_{p,q}$ et donc, quand on rajoute des ps.\ droites, jamais bigonal.
\marginal{cas mutuellement exclusifs}
\note{les trois cas « On sait \ldots » sont réunis à gauche par une accolade, qui porte la note marginale}

On sait $\exists\, K \simeq \mathrm{III}_6$ \uncertain{ou} $\mathrm{III}_7$, donc $\exists\, K_6 \simeq \mathrm{III}_6$, et \struck{\ill{}}, \ill{} $\mathrm{III}_7$, si on rajoute une ps.\ droite, elle n'est pas bigonale.

Le \uncertain{cas} de départ \struck{\ill{}} \ill{} \ill{} $\mathrm{II}_4$ ou $\mathrm{III}_6$.
\note{en bas à droite, un dessin : cinq droites, dont deux points de concours marqués}

\section{Arrangements de 3 pseudo-droites}
\note{titre de sa main, en tête de la page 31, précédé de « I) »}

\page{31}
\note{numéroté \emph{1} (entouré) en haut à gauche, de sa main ; les feuillets suivants portent 2 à 5 aux pages 33, 35, 37 et 39 : il ne numérote que les rectos, les pages paires en sont le verso et continuent le texte. Encre brune, écriture posée}
I) Arrangements de $3$ pseudo-droites

1) \emph{Cube combinatoire} (de dim $3$). La donnée équivaut à celle d'une famille $(\Phi_i)_{i \in I}$ de trois ens.\ à deux éléments ($\operatorname{card} I = 3$, $\operatorname{card} \Phi_i = 2$ $\forall i \in I$)
\[
(1) \qquad
\begin{array}{cl}
\Phi = \coprod_{i \in I} \Phi_i & \operatorname{card} 6 \\
\downarrow & \\
I & \operatorname{card} 3
\end{array}
\]
Les éléments de $\Phi$ s'identifient aux faces du cube, ceux de $I$ aux paires de faces opposées. Les \uncertain{facettes} du cube s'identifient aux \emph{sections partielles} \add{des parties $\neq \emptyset$ de $I$} de $\Phi$ sur $I$, avec la relation d'incidence correspondant à la relation \emph{opposée} \struck{au} à celle du prolongement des sections partielles, i.e.\ de l'inclusion des parties correspondantes de $\Phi$ (parties $A \subset \Phi$ non vides, telles que l'appl.\ induite $A \to I$ soit injective).
\[
(2) \qquad
\begin{array}{l}
C_2 \simeq \Phi \simeq \mathfrak{P}_1(\Phi) \\
C_1 \simeq \lbrace A \in \mathfrak{P}_2(\Phi) \mid A \to I \text{ injectif} \rbrace \\
C_0 \simeq \lbrace A \in \mathfrak{P}_3(\Phi) \mid A \to I \text{ injectif (donc bij.)} \rbrace
\end{array}
\]
\note{« des parties $\neq \emptyset$ de $I$ » est écrit au-dessus de la ligne. $\mathfrak{P}$ rend son P gothique (ensemble des parties) ; « donc » est lu sans certitude}

On a un automorphisme \add{(involutif)} canonique $\underline{a}$ de (1), induisant l'identité dans $I$ et la transposition dans chacun des $\Phi_i$, qui induit un automorphisme \add{(central)} du cube (2), lequel est l'antipodisme.

Si on prend le polyèdre opposé au cube, on trouve l'octaèdre donné par
\[
(3) \qquad
\left\lbrace
\begin{array}{l}
O_0 = C_2 \simeq \Phi \simeq \mathfrak{P}_1(\Phi) \\
O_1 = C_1 \simeq \lbrace A \in \mathfrak{P}_2(\Phi) \mid A \hookrightarrow I \rbrace \\
O_2 = C_0 \simeq \lbrace A \in \mathfrak{P}_3(\Phi) \mid A \hookrightarrow I \rbrace
\end{array}
\right.
\]
avec cette fois-ci la relation \add{d'incidence qui est celle} \struck{de} du prolongement des sections

\page{32}
partielles i.e.\ l'inclusion des parties de $\Phi$.

Passant au quotient par $\underline{a}$, on trouve le \emph{cube combinatoire} \add{\emph{gauche} (ou} \emph{projectif}) $\widetilde{C}$, et l'\emph{octaèdre comb.} \add{\emph{gauche} (ou} \emph{projectif}) $\widetilde{O}$, associés à $\underline{a}$,
\[
(4) \qquad \widetilde{C}_i = C_i/\lbrace 1, \underline{a} \rbrace
\qquad \widetilde{O}_i = O_i/\lbrace 1, \underline{a} \rbrace
\]
(relation d'incidence \uncertain{déduite} par passage au quotient). L'octaèdre \struck{projectif} \add{gauche} a donc $3$ sommets, $6$ arêtes et $4$ faces --- en fait c'est la décomposition du plan projectif réel par un système de trois droites indexées par $I$. (À tout élément $i \in I$, on fait correspondre la « droite » correspondant aux deux « arêtes » projectives, qui correspondent aux paires \uncertain{d'arêtes} ordinaires données par $A \in \mathfrak{P}_2(\Phi)$, $A \xrightarrow{\sim} I \setminus \lbrace i \rbrace$.)
\note{« gauche (ou » est ajouté en interligne devant « projectif », souligné comme lui ; plus bas « gauche » remplace « projectif » biffé. Le mot se lit sûrement « gauches » à la page 37 (« oct.\ gauches »). À gauche, un dessin : trois droites en position générale formant un triangle}

On en récupère le cube ou octaèdre ordinaire, à isomorphismes \uncertain{canoniques} près, à partir de la figure « gauche » \uncertain{projective} correspondante, un \ill{} \ill{}, dans le cas ordinaire, \ill{} que les groupes d'\uncertain{automorphismes}, dans le cas \struck{projectif} \add{« gauche »}, \ill{} \ill{} par \ill{}, plus précisément, $C \mapsto \widetilde{C}$ resp.\ $O \mapsto \widetilde{O}$
\note{les deux dernières phrases sont d'une lecture incertaine ; « isomorphismes » et « gauche » y sont ajoutés en interligne}

\page{33}
\note{sa page \emph{2}, en haut à gauche}
ne sont pas pleinement fidèles mais \uncertain{sont} fidèles, car l'antipodisme $\underline{a}$ est un automorphisme non trivial de $C$, $O$ et induit l'automorphisme trivial de $\widetilde{C}$, $\widetilde{O}$. Mais on trouve une équivalence de catégories en prenant des cubes comb.\ \emph{orientés}, et des octaèdres combinatoires \emph{orientés}.
\[
(5) \qquad
\begin{array}{ccc}
\text{cubes comb.\ orientés} & \xrightarrow{\ \approx\ } & \text{cubes comb.\ gauches} \\
\wr\downarrow & & \wr\downarrow \\
\text{octaèdres comb.\ orientés} & \xrightarrow{\ \approx\ } & \text{octaèdres comb.\ gauches}
\end{array}
\]
\note{dans les deux lignes de (5), « projectifs » est biffé et remplacé par « gauches » ; la flèche verticale de droite n'est marquée que d'un $\wr$}

L'ens.\ des orientations du cube, ou de l'octaèdre, donné par (1), est \struck{\ill{} \ill{}} \add{donné à \uncertain{isom.}\ près} \ill{}
\[
(6) \qquad \operatorname{or}(C) \simeq \operatorname{or}(O) \simeq \Bigl(\bigwedge_{i \in I} \Phi_i\Bigr) \wedge \operatorname{or}(I)
\]
où $\operatorname{or}(I)$ est l'ens.\ des deux ordres circulaires sur $I$ (ou des deux « orientations » de l'ens.\ fini $I$) ; à partir de $\widetilde{C}$ ou $\widetilde{O}$, on récupère $C$ ou $O$ \struck{orienté comme un \ill{} orientation} comme le revêtement orienté canonique.
\marginal{la donnée d'une orientation de $C$, ou $O$, \struck{\ill{}} revient à la donnée d'un iso (6 bis) $\bigl(\bigwedge_{i \in I} \Phi_i\bigr) \wedge \operatorname{or}(I) \simeq \underline{1}$ ou encore $\bigwedge_{i \in I} \Phi_i \simeq \operatorname{or}(I)$}
\note{la note marginale est écrite verticalement le long du bord gauche, en face de (6) ; $\bigwedge$ est son signe en chapeau, le produit « extérieur » de deux ensembles à deux éléments ; $\underline{1}$ est souligné}

Soit $i \in I$, d'où « droite » $D_i$ de la figure $\widetilde{O}$. Les deux arêtes de $D_i$ correspondent aux deux bijections $\Phi_j \simeq \Phi_k$ (où $I = \lbrace i, j, k \rbrace$), ou \uncertain{intrinsèquement}, aux éléments de
\[
(7) \qquad A_i = \bigwedge_{\alpha \in I \setminus \lbrace i \rbrace} \Phi_\alpha
\qquad \text{(\uncertain{ens.}\ des deux arêtes de $\widetilde{O}$ portées par $D_i$)}
\]
\struck{Les orient}
L'un des deux sommets dans $D_i$, d'autre part,

\page{34}
est donné par
\[
(8) \qquad S_i \simeq I \setminus \lbrace i \rbrace
\qquad \text{(ens.\ des \add{deux} sommets sur la droite)}
\]
L'ens.\ des deux orientations de $D_i$ (considéré comme ensemble d'un bigône) est donc
\[
(9) \qquad \Omega_i \simeq S_i \wedge A_i \simeq (I \setminus \lbrace i \rbrace) \wedge \underbrace{\bigwedge_{\alpha \in I \setminus \lbrace i \rbrace} \Phi_\alpha}_{\Phi_i \wedge \operatorname{or}(I) \text{ à cause de (6 bis)}}
\]
\note{la lettre $S_i$ rend son S cursif (sommets) ; « deux » est ajouté au-dessus de la ligne. L'accolade sous le dernier facteur de (9) porte « $\Phi_i \wedge \operatorname{or}(I)$ à cause de (6 bis) »}

Notons que
\[
(10) \qquad \underbrace{\bigwedge_{i \in I} A_i}_{\bigwedge_{\substack{\alpha, i \in I \\ \alpha \neq i}} \Phi_\alpha \;=\; \bigwedge_{\alpha \in I} \Phi_\alpha^2 \;\simeq\; \underline{1}} \simeq \underline{1}
\]
\note{au-dessus du membre de droite de (10), une seconde accolade porte « $\Phi_i \wedge \bigwedge_{\alpha \in I \setminus \lbrace i \rbrace} (I \setminus \lbrace \alpha \rbrace)$ canoniquement », d'une lecture incertaine ; $\underline{1}$ est son 1 à double jambage (l'ensemble à un élément)}

d'où
\[
(11) \qquad \bigwedge_{i \in I} \Omega_i \overset{\varphi}{\simeq} \bigwedge_{i \in I} S_i
\quad \Bigl[\simeq \bigwedge_{i \in I} (I \setminus \lbrace i \rbrace) \simeq \operatorname{or}(I)\Bigr]
\]
\note{sous le dernier $\simeq$ de (11) : « isom.\ can.\ évident associé à l'ens.\ à $3$ éléments $I$ »}

\emph{NB} Quand on part d'une famille, i.e.\ de trois bigônes $(S_i, A_i)_{i \in I}$, la donnée d'un isom.\ (10) équivaut à celle d'un isom.\ $\varphi$ (11) (qui, pour $S_i = I \setminus \lbrace i \rbrace$ $\forall i \in I$), se réduit à un isom.
\[
(11\ \text{bis}) \qquad \bigwedge_{i \in I} \Omega_i \simeq \operatorname{or}(I)
\]
ou encore à un isom.
\[
(12) \qquad \Bigl(\bigwedge_{i \in I} \Omega_i\Bigr) \wedge \operatorname{or}(I) \simeq \underline{1}
\]
\note{le mot après « NB » (lu « Quand ») est surchargé ; « i.e.\ » est biffé puis récrit}

\page{35}
\note{sa page \emph{3}, en haut à gauche}
qui peut s'interpréter aussi de la façon suivante : la famille \add{$(\Omega_i)_{i \in I}$} des trois ensembles $\Omega_i$ à deux éléments définit un cube --- et \uncertain{on} a une orientation de ce cube.

\emph{NB} Les $\Omega_i$ et $\Phi_i$ se déterminent mutuellement par
\[
(13) \qquad
\left\lbrace
\begin{array}{l}
\Omega_i \simeq \Phi_i \wedge \bigwedge_{\alpha \in I \setminus \lbrace i \rbrace} (I \setminus \lbrace \alpha \rbrace) \\[1ex]
\Phi_i \simeq \Omega_i \wedge \bigwedge_{\alpha \in I \setminus \lbrace i \rbrace} (I \setminus \lbrace \alpha \rbrace)
\end{array}
\right.
\]
qui amènent des isom.\ canoniques
\[
(14) \qquad
\left\lbrace
\begin{array}{l}
\bigwedge_{i \in I} \Omega_i \simeq \bigwedge_{i \in I} \Phi_i \\[1ex]
\underbrace{\Bigl(\bigwedge_{i \in I} \Omega_i\Bigr) \wedge \operatorname{or} I}_{\text{orientations de l'oct.\ comb.\ défini par } (\Omega_i)_{i \in I}}
\simeq
\underbrace{\Bigl(\bigwedge_{i \in I} \Phi_i\Bigr) \wedge \operatorname{or} I}_{\text{orientations de l'oct.\ comb.\ défini par } (\Phi_i)_{i \in I}}
\end{array}
\right.
\]
\marginal{\ill{} $A_i$, $\Phi_i$, $\Omega_i$, puisque \ldots\ $A_i = \bigwedge (I \setminus \lbrace i \rbrace) \ldots \simeq \Phi_i \wedge \bigwedge_{\alpha \in I} \ldots \operatorname{or}(I)$ \ldots\ $\simeq \Phi_i$ (cf.\ (9))}
\note{la note marginale est écrite en diagonale dans le coin gauche, en face de (13), avec plusieurs mots noircis ; on n'en donne que les formules lisibles}

En résumé

\emph{Proposition} Considérons la carte combinatoire \add{\ill{}} (avec trois sommets, six arêtes, quatre faces qui sont des triangles) correspondant à un « arrangement » (combinatoire) de \struck{quatre} \add{trois} ps.-droites \add{$D_i$ ($i \in I$)} dans \add{\uncertain{un}} plan projectif ». Soit $K$ son
\marginal{octaèdre gauche}
\note{la parenthèse et la note marginale « octaèdre gauche » désignent la même carte. Un petit signe entouré, après « combinatoire », renvoie à un ajout que l'on ne lit pas}

\page{36}
squelette de dimension $1$ » ($1$-complexe combinatoire, muni des systèmes des trois \struck{ps.}\ sous-$1$-complexes $D_i$. \struck{\ill{}} Donc $K$ est la réunion des sous-$1$-complexes $D_i$, qui sont des bigônes, \add{chacun} défini\supplied{s} par \add{$S_i$} \add{\uncertain{un}} ens.\ de \add{deux} sommets et d'arêtes \add{$A_i$} \struck{\ill{}} :
\note{« chacun », « $S_i$ », « $A_i$ », « deux » sont des ajouts interlinéaires, reliés par des traits}
\struck{\ill{}}
\[
S_i \simeq I - \lbrace i \rbrace
\]
\note{avant la formule, une accolade ouvre deux lignes : la première, « $S_i \simeq I \setminus \lbrace i \rbrace$, $A_i$ », est barrée de deux traits ; la seconde la récrit. Un signe biffé précède l'accolade}
(\emph{NB} \struck{\ill{}} l'\uncertain{ens.}\ des sommets de $K$ (i.e.\ de \uncertain{$\widetilde{C}$}) est en corr.\ biun.\ avec $I$ \uncertain{en associant} à chaque $i \in I$ l'unique sommet qui n'est pas situé sur $D_i$). Soit
\[
\Omega_i = S_i \wedge A_i = (I - \lbrace i \rbrace) \wedge A_i
\]
l'ens.\ des orientations du bigône $D_i$, i.e.\ de la ps.\ droite $D_i$. Soit
\[
\begin{aligned}
\Phi_i &= \Omega_i \wedge \bigwedge_{\alpha \in I \setminus \lbrace i \rbrace} (I \setminus \lbrace \alpha \rbrace) \\
&\simeq A_i \wedge \bigwedge_{\alpha \in I} (I \setminus \lbrace \alpha \rbrace)
\end{aligned}
\]
de sorte que l'on a des isom.\ canoniques
\[
(15) \qquad \bigwedge_{i \in I} \Omega_i \simeq \bigwedge_{i \in I} \Phi_i \simeq \Bigl(\bigwedge_{i \in I} A_i\Bigr) \wedge \underbrace{\bigwedge_{\alpha \in I} (I \setminus \lbrace \alpha \rbrace)}_{\operatorname{or} I}
\]
les ens.\ $S_i$, $A_i$, $\Omega_i$, $\Phi_i$ $(i \in I)$ et les isom.
\note{la phrase qui commence par « (NB » ouvre deux parenthèses et n'en ferme qu'une avant « Soit » ; transcrite telle quelle. Le premier mot du NB est surchargé}

\page{37}
\note{sa page \emph{4}, en haut à gauche}
(15) étant définis donc directement en termes du $1$-complexe \add{comb.}\ $K$ et de la donnée des trois ps.\ droites $D_i \subset K$ dans $K$. On peut réécrire (15) comme une \uncertain{suppl.}\ d'iso
\[
(15\ \text{bis}) \qquad
\underbrace{\Bigl(\bigwedge_{i \in I} \Omega_i\Bigr) \wedge \operatorname{or} I}_{\substack{\text{or.\ de l'oct.} \\ \text{défini par } (\Omega_i)_{i \in I}}}
\simeq
\underbrace{\Bigl(\bigwedge_{i \in I} \Phi_i\Bigr) \wedge \operatorname{or}(I)}_{\substack{\text{or.\ de l'oct.} \\ \text{défini par } (\Phi_i)_{i \in I}}}
\simeq \bigwedge_{i \in I} A_i
\]

Ceci posé, il y a, \struck{un isom.\ canonique} \add{du fait que} \emph{$K$ provient de la carte $C$}, \struck{un iso} \add{\emph{trivialisation}} \emph{canonique} \emph{des} $\mathbb{Z}/2\mathbb{Z}$-\emph{torseurs} (15 bis). De façon précise, le foncteur
\[
(16) \qquad (\text{oct.\ gauches}) \longrightarrow \overbrace{\text{graphes octaédraux gauches}}^{\text{gr.\ oct.\ gauches}^{+}}
\]
munis d'une trivialisation $\bigwedge_{i \in I} A_i \simeq \underline{1}$ (où $I$ est l'ens.\ des $3$ sommets, et $A_i$ est l'ens.\ des $2$ arêtes qui ne passent pas par le sommet $i$) \marginal{($\simeq$ gr.\ oct.\ \uncertain{gauches} \ill{} avec structure suppl.\ une orientation des gr.\ octaédraux gauches)} est une \emph{équivalence} de catégories.
\note{la note marginale, en petite écriture à gauche de l'accolade, est d'une lecture incertaine ; « 3 » et « 2 » sont ajoutés en interligne}

On peut résumer les relations entre octaèdres et graphes octaédraux, gauches ou ordinaires, orientés ou non, \ill{} \add{dans le} diagramme commutatif de foncteurs suivant \add{(\uncertain{notation} \uncertain{abrégée} \ill{})} :

(17)

\begin{tikzcd}[column sep=small, row sep=small, nodes={font=\scriptsize}]
\text{oct}^{+} \arrow[rr, "\approx"] \arrow[dr, "\approx"] \arrow[dd] & & \text{oct.\ gauches} \arrow[dr, "\approx"] \arrow[dd, no head, "=" description] & \\
& \text{gr.\ oct}^{+} \arrow[rr, "\approx"] \arrow[dd] & & \text{gr.\ oct.\ gauches}^{+} \arrow[dd] \\
\text{oct} \arrow[rr] \arrow[dr, "\approx"] & & \text{oct.\ gauches} \arrow[dr] & \\
& \text{gr.\ oct} \arrow[rr, "\approx"] & & \text{gr.\ oct.\ gauches}
\end{tikzcd}

\note{diagramme en cube : les deux faces, arrière (octaèdres) et avant (graphes octaédraux), sont reliées par les obliques. Le trait vertical de droite à l'arrière est un double trait « $\|$ » (égalité) interrompu par la flèche horizontale qu'il croise ; les quatre autres verticales sont des flèches, trois d'entre elles marquées à leur origine d'un petit crochet surmonté de points, que l'on ne sait pas interpréter. La flèche horizontale du milieu à l'arrière ne porte pas de $\approx$. Une virgule suit le diagramme}

\page{38}
où les flèches marquées $\approx$ sont des équivalences, les flèches marquées $\hookrightarrow$ sont \struck{\ill{}} fidèles, les flèches $\twoheadrightarrow$ sont \struck{\ill{}} surjectives sur les ens.\ d'isomorphismes, la flèche \emph{verticale} $\|\!\downarrow$ est l'identité de (oct.\ g.). La face supérieure est formée des objets « orientés », les quatre flèches constituant ce carré sont des équivalences. La face inférieure \struck{\ill{}} \add{formée} des objets \uncertain{non} orientés, \struck{\ill{}} les deux flèches
\[
\text{oct} \longrightarrow \text{gr.\ oct.} \longrightarrow \text{gr.\ oct g.}
\]
sont \struck{des} \add{des} équivalences, les deux autres sont resp.\ surjectives sur \add{les} \uncertain{isom.}, et fidèles. Les quatre flèches de passage de la théorie orientée vers la non orientée sont fidèles. La face de gauche est formée des objets « ordinaires » \struck{\ill{}} « sphériques », celle de droite des objets « gauches » ou « projectifs », les \add{quatre} flèches de passage de l'une à l'autre sont des équivalences, à l'exception de
\[
\text{oct} \longrightarrow \text{oct.\ gauches} \qquad (\text{surjectif sur les isom.}).
\]
La face postérieure du cube est formée des configurations deux-dimensionnelles, \struck{celle} celle de la face antérieure est formée des configurations $1$-dimensionnelles associées, les quatre flèches de passage sont des équivalences, sauf
\[
\text{oct.\ g} \longrightarrow \text{gr.\ oct.\ gauches},
\]
qui est fidèle et fait l'objet de l'énoncé (16) plus haut.
\marginal{correspond.\ \struck{\ill{}} $\mathfrak{S}_4 \times \lbrace \pm 1 \rbrace$ \ldots\ $\mathfrak{S}_4$}
\note{les deux premières sortes de flèches sont dessinées dans le texte, l'une avec un petit crochet à l'origine, l'autre avec un crochet à la pointe ; on les rend par $\hookrightarrow$ et $\twoheadrightarrow$. La note marginale, en diagonale dans le coin supérieur gauche, relie par deux flèches les groupes $\mathfrak{S}_4 \times \lbrace \pm 1 \rbrace$ et $\mathfrak{S}_4$ (lecture incertaine) aux deux premières lignes ; plusieurs signes y sont noircis. Les « faces » sont celles du cube (17) de la page 37}

\page{39}
\note{sa page \emph{5}, en haut à gauche}
On appellera \emph{orientation} de $(K, (D_i)_{i \in J})$, la donnée, pour $\forall\, I \in \mathfrak{P}_3(J)$ comme \uncertain{dessus} (\uncertain{de} l'ens.\ de $3$ pseudo-droites non concourantes), d'une \emph{orientation} \add{$\omega_I$} du graphe octaédral gauche correspondant.
\note{$\mathfrak{P}_3(J)$ : l'ensemble des parties à trois éléments de $J$ ; « $\omega_I$ » est ajouté en interligne}

Soit maintenant $\mathcal{X}$ une \emph{carte} qui soit un arrangement de pseudo-droites, i.e.\ son $1$-squelette \add{$K$} est donc muni d'une famille $(D_i)_{i \in J}$ de pseudo-droites, i.e.\ \struck{\ill{}} de un arrangement $1$-dim.\ de ps.-droites. Il plus, il est muni d'une orientation
\[
\omega = (\omega_I)_{\substack{I \in \mathfrak{P}_3(J) \\ \bigcap_{i \in I} D_i = \emptyset}}
\]
On trouve ainsi un foncteur
\[
(18) \qquad \mathcal{X} \longmapsto (K, (D_i)_{i \in J}, \omega)
\]
de la catégorie des arrangements comb.\ de pseudo-droites, vers la catégorie combinatoire des arrangements $1$-dimensionnels orientés de pseudo-droites. \struck{\ill{}}
\note{le « $1$ » de « $1$-dimensionnels » est ajouté en interligne ; au-dessous, un mot commencé et laissé en suspens}

\emph{Th.} Le foncteur précédent est \emph{pleinement fidèle}, i.e.\ pour deux arrangements comb.\ $\mathcal{X}$, $\mathcal{X}'$ de ps.\ droites, on a
\[
\operatorname{Isom}(\mathcal{X}, \mathcal{X}') \xrightarrow{\ \sim\ } \operatorname{Isom}\bigl((K, (D_i)_{i \in J}, \omega), (K', (D'_{i'})_{i' \in J'}, \omega')\bigr)
\]
\note{« Th. » est surchargé et souligné. Dans la marge gauche, en face du théorème, trois lignes obliques de sa main sont barrées : on y devine « Il \ldots\ \ldots\ pour \ldots\ \ldots\ \ldots\ qui \ldots\ \ldots\ bigone », lecture trop incertaine pour être donnée. Un trait oblique, à droite, prolonge la dernière formule}

\page{40}
Soit $K$ un $1$-complexe combinatoire fini, $(D_i)_{i \in J}$ une fam.\ finie de sous-complexes polygonaux, tels que toute arête de $K$ appartienne à un $D_i$ et un seul, \struck{et que $K$ \ill{} \ill{} i.e.\ $\forall$ sommet de $K$ \ill{} \ill{}} \struck{\ill{} \ill{} au moins un des $D_i$}. On suppose de plus que pour $i \neq j$, $D_i$ et $D_j$ se \uncertain{rencontrent} en un sommet et un seul, et que $\forall$ sommet \add{de $K$} \uncertain{appartient} à au moins $2$ des $D_i$. On dit alors que $(K, (D_i)_{i \in J})$ est un \add{\struck{\ill{}} $1$-dimensionnel} arrangement de ps.\ droites. Supposons que l'intersection des $D_i$ soit vide, \add{i.e.\ que} $K$ ne soit pas \uncertain{réduit} à \uncertain{un} « point » \add{(ou « boucles »)} \struck{monogones} \struck{polygones} \uncertain{constitués} par \ill{} \ill{}. Il existe alors une partie $I \subset J$ de cardinal \ill{}. \struck{\ill{}}

Pour \uncertain{une} partie $I$ de $J$, soit $K_I$ le sous-complexe « topologique » \add{réunion} des $D_i$ ($i \in I$) \uncertain{et} des sommets \uncertain{que les} intersections \ill{} \ill{}. \uncertain{Si} (card $I = 3$, $\bigcap_{i \in I} D_i = \emptyset$), $K_I$ est un graphe octogonal gauche, donc on peut parler d'\emph{orientation} de ce graphe octogonal g.
\marginal{\emph{NB} L'ordre \ill{} d'un \ill{} \ill{} \ill{} \ill{} bigônes \ill{} $\geq 2$ \ldots\ $K_I$ \ldots\ $(D_i)_{i \in I}$ \ldots}
\note{page très reprise, écrite vite : plusieurs passages biffés d'un trait se superposent aux ajouts interlinéaires, et seule la charpente des phrases se lit. La note marginale, en longues lignes obliques dans le coin inférieur gauche, n'est lisible que par fragments. « octogonal » (deux fois) est transcrit tel qu'il se lit ; le contexte, et la page 37 (« graphes octaédraux gauches »), attendent « octaédral »}

\end{document}
